\section{口服药物吸收的评价方法}

\begin{frame}
    \frametitle{口服药物吸收的评价方法}
    口服药物吸收评价主要通过药物渗透性进行。药物渗透性的测定可采用人体试验方法或其他能预测药物体内吸收程度的非人体试验方法。
    \begin{enumerate}
      \item 体内法（in vivo）
      \item 在体法（in situ）
      \item 体外法（in vitro）
    \end{enumerate}
\end{frame}


\begin{frame}
    \frametitle{体内法}
    体内法是以整体动物机体为研究对象，进行药动学研究。

    口服给药后，在不同时间点采集生物样本如血液、尿液，再通过一定方法测定其药物浓度，计算达峰时间（$T _\mathrm{max}$）、达峰浓度（$C_\mathrm{max}$）、药时曲线下面积（$AUC$）等药动学参数，进而评价药物吸收的程度和速度。

    \begin{itemize}
      \item 给药方法一般是大动物口服给药，小动物灌胃给药。
      \item 此法采用整体动物进行实验，能真实地反映药物口服后的体内吸收情况，但是它综合了物化、生理、剂型等众多因素的结果，不能特异性地反映肠道对药物的吸收作用。
      \item 由于实验周期长、操作相对繁琐、影响因素较多、动物个体差异较大等原因，故体内法很少应用于药物吸收机制的研究，一般用于研究药物体内药动学特征。
    \end{itemize}

\end{frame}


\begin{frame}
    \frametitle{在体法}      
    在体回肠灌流法：在动物麻醉状态下，不切断血管和神经，将待研究的肠段两端结扎，肠管插管，洗净肠内容物，构成循环回路，使药物在肠腔内循环。不同时间测定灌流液中药物浓度变化，可以获得药物在肠内吸收情况。本法能避免胃容物和消化道固有生理活动对结果影响。
      \begin{itemize}
        \item 在体法是建立在整体动物水平上的实验，与体内法不同的是其将干扰因素大大减少，同时保证了肠道神经与内分泌调节的完整性，也保证了血液与淋巴液供应，使得待考察组织的生物活性大大提升，能够较为准确地反映药物在肠道的真实吸收情况，常用于研究药物的渗透和吸收动力学。
        \item 在体法由于是建立在整体水平之上，因此个体差异较大，对实验动物的数量有一定要求，以保证最小的数据变异。
      \end{itemize}

\end{frame}


\begin{frame}
    \frametitle{体外法}
    \begin{enumerate}
        \item 肠外翻囊法
        
        动物麻醉后， 取出一段肠管，用pH 7.40 Krebs-Henseleit (K-H)生理液冲洗，将肠管套于玻璃棒上外翻， 一端固于聚乙烯管上， 另一端结扎，内注入一定量的pH 7.40 K-H生理溶液，置于含有药物的pH 7.40 K-H生理溶液中，37℃恒温， 在5\%O$_2$气流下，定时从内外两侧取样，测定药物浓度变化。该方法可根据需要研究不同肠段的药物吸收或分泌特性及其影响因素。
        \item 细胞培养模型法
      \end{enumerate}
      \begin{itemize}
        \item 操作简单、试验周期短，可以用于药物早期高通量筛选。
        \item 肠道的不同节段以及缺乏血液、淋巴液的供应对细胞旁路通道和酶活性造成的影响均会影响实验结果，其组织活性与其在体内环境中会存在较大的差异，不能反映药物在肠道环境中的真实吸收情况。
      \end{itemize}
\end{frame}


\begin{frame}
    \frametitle{Caco-2细胞模型}
    细胞模型是体外法评价药物吸收的重要手段，最常用的是Caco-2细胞模型。

    Caco-2细胞模型(human colon adenocarcinoma
cell line ) 最早是由 美国 堪萨斯大学 药学 院
Ronald T. Borchardt 博士的实验室以及英国 的
Ciba-Geigy实验室在二十世纪八十年代后期建立的。
是第一个被国外制药企业以及实验室广泛应用的肠
道吸收的细胞模型，也是迄今为止的最好的肠道转
运模型和上皮转运模型。

Caco-2细胞系来源于人的直肠癌， 其结构和生
化作用类似于人小肠上皮细胞， 含有与小肠刷状缘
上皮相关的酶系， 与正常的成熟小肠上皮细胞在体
外培育过程中出现逆向分化(反分化) 不同， Caco-2
细胞在传统的细胞培养条件下， 生长在多孔的可渗
透的聚碳酸酯膜上可达到融合并自发分化为肠上皮
细胞， 形成连续的单层(monolayer) 。培育约3星期后， 其电阻值达到400 Ω （有的
要求达到600Ω ） ， 它们形成致密的细胞单层，
结构和生化作用类似于小肠上皮细胞， 3-5天
内适合进行转运试验。
\end{frame}


% \begin{frame}
%     \frametitle{Caco-2细胞模型}
% Caco-2细胞单层膜模型构建成功后，在膜的两侧分为了A侧（肠腔侧）与B侧（基底侧）。

% \begin{figure}[htpb]
%     \centering
%     \includegraphics[width=.8\linewidth]{caco2}
% \end{figure}
% \end{frame}


\begin{frame}
    \frametitle{Caco-2细胞模型的优点}
    \begin{enumerate}
      \item Caco-2细胞来源是人结肠癌细胞，同源性好。该模型与细胞在肠中的吸收有良好的相关性、较高的重现性；
      \item 在细胞水平上研究药物在小肠粘膜中的吸收、转运和代谢
      \item 与动物实验相比， Caco-2细胞生命力强， 培养方法简单，更经济高效；
      \item 可同时对大批量药物进行快速筛选， 所需药物量少， 仅需要数毫克即可；
      \item 可以从不同方向进行摄取和转运实验，可用于区分肠腔内不同吸收途径的差别；
      \item 有较广泛的适用性（可用于原料药、制剂、大小分子的跨肠膜研究）
      \item 可以同时研究药物对粘膜的毒性
    \end{enumerate}
\end{frame}


\begin{frame}
    \frametitle{Caco-2细胞模型的局限}
    \begin{enumerate}
      \item 缺少粘蛋白产生细胞， 缺少肠壁粘液层；
      \item 具有较高的细胞通透性， 若喂养次数多， 某些酶载体或细胞内容物溢出；
      \item 缺少某些代谢酶；
      \item 细胞的培养和代系对于Caco-2细胞形态学和生理性质有影响， 引起细胞对药物转运产生差别；
      \item 由于培养条件差别， 使不同实验室出来结果缺乏可比性
    \end{enumerate}
\end{frame}


\section{药物在其它部位的吸收}


\begin{frame}
    \frametitle{在口腔粘膜中吸收}
    药物经口腔黏膜给药可发挥局部或全身治疗作用，局部作用的剂型多为溶液型或混悬型，如漱口剂 、气雾剂 、膜剂、软膏剂、口腔片剂等 。 欲产生全身作用常用舌下片 、黏附片、贴剂等 。产生全身作用的口腔用药主要部位是颊黏膜，其次是舌下黏膜 。

    \begin{figure}[htpb]
        \centering
        \includegraphics[width=.5\linewidth]{口腔给药}
    \end{figure}
\end{frame}

\begin{frame}
    \frametitle{在口腔粘膜中吸收}
    \hspace{-2em}
    \begin{minipage}{0.28\linewidth}
        \begin{figure}[htpb]
            \centering
            \includegraphics[width=\linewidth]{颊粘膜}
            \caption{颊粘膜结构示意图}
        \end{figure}
    \end{minipage}\hfill
    \begin{minipage}{0.68\linewidth}
         颊黏膜可分为四层结构 ： 
        \begin{enumerate}
          \item 上皮层：由角质形成细胞与非角质形成细胞组成 ， 为药物透过黏膜的主要屏障 。
          \item 基底层
          \item 固有层：致密的结缔组织
          \item 黏膜下层：疏松的结缔组织
        \end{enumerate}
        药物在口腔黏膜中渗透到达黏膜下层 ， 黏膜下层有大量的毛细血管汇总至颈内静脉而进入血液循环 ， 因此可避免肝脏首过效应 。
    \end{minipage}
\end{frame}

\begin{frame}
    \frametitle{在直肠粘膜中吸收}
    直肠的吸收面积不大，但血流量较为丰富，药物容易吸收。直肠给药可用于局部治疗或发挥全身作用 ， 常用的剂型是栓剂或灌肠剂 。 栓剂用于全身治疗有许多优点 ： 
    \begin{enumerate}
      \item 药物直肠吸收后 ， 大部分绕过肝脏进入大循环发挥全身作用 ， 降低了肝首过效应 ， 也相应减少药物对肝脏的毒副作用； 
      \item 避免胃肠 pH 和酶的影响和破坏 ， 避免药物对胃肠功能干扰 ， 对胃有刺激的药物可采用直肠给药 ；
      \item 作用的时间一般比片剂长 ， 通常给药为 1 、 2 次 / 日；
      \item 适于不愿或不能吞服药物的患者 ， 对于婴幼儿及神志障碍的患者 ， 使用栓剂较口服或注射给药更容易 、 更安全 ；
      \item 可作为多肽蛋白质类药物的吸收部位 。
    \end{enumerate}
\end{frame}

\begin{frame}
    \frametitle{在直肠粘膜中吸收}
    直肠给药制剂首先需要药物从基质中释放，释放的药物穿过上皮细胞进入血管被吸收。
    \begin{figure}[htpb]
        \centering
        \subfloat[药物从基质中的释放与吸收过程]{\includegraphics[height=5cm]{直肠给药}}\quad
        \subfloat[直肠与肛门部位的血管分布]{\includegraphics[height=5cm]{直肠血管分布}}
    \end{figure}
\end{frame}

\begin{frame}
    \frametitle{透皮肤吸收}
    皮肤给药可以用于局部皮肤病的治疗，也可以经皮肤吸收后治疗全身性疾病。对于皮肤病，由于病灶部位的深浅不同 ，某些药物需要透过角质层以后才能起效；而对于全身性疾病 ， 药物必须通过角质层 ，被皮下毛细血管吸收进入血液循环后才能起效。经皮吸收是指药物从应用于皮肤上的制剂中释放与穿透皮肤进入体循环的过程 。

    \hspace{-2em}
    \begin{minipage}{0.38\linewidth}
      \begin{figure}[htpb]
        \centering
        \includegraphics[width=\linewidth]{皮肤结构}
      \end{figure}
    \end{minipage}\hfill
    \begin{minipage}{0.58\linewidth}
      \begin{itemize}
        \item 角质层是经皮吸收药物的主要屏障。
        \item 当皮肤角质层受损时，药物的通透性显著增加。
        \item 不同部位的角质层厚度不同，不同的药物在不同的部位吸收程度是不同的。
        \item 皮肤有较大的表面积, 皮肤吸收往往也是毒物（如农药）等吸收造成中毒的重要途径之一。 
      \end{itemize}
    \end{minipage}
\end{frame}


\begin{frame}
    \frametitle{在肺部中吸收}
    肺部给药又称吸入给药， 主要是通过口腔吸入 ， 经过咽喉进入呼吸道 ， 到达呼吸道深处或肺部 ， 起到局部作用或吸收后产生全身治疗作用 。 

    \begin{itemize}
      \item 气体或挥发性药物可以在肺上皮细胞或气管粘膜吸收
      \item 用于肺部给药的剂型包括气雾剂 、 喷雾剂和粉雾剂 
      \item 肺吸收也往往是毒品和环境毒物重要的吸收途径
      \item 优点是药物吸收快，可避免肝肠分解
      \item 缺点是药物的剂量难以控制，药物也会对肺上皮细胞产生刺激
    \end{itemize}
\end{frame}


\begin{frame}
  \frametitle{在肺部中吸收}
  \hspace{-2em}
  \begin{minipage}{0.4\linewidth}
    \begin{figure}[htpb]
      \centering
      \includegraphics[width=\linewidth]{呼吸系统}
      \caption{人体呼吸器官示意图}
    \end{figure}
  \end{minipage}\hfill
  \begin{minipage}{0.58\linewidth}
    \begin{itemize}
      \item 气体进入口腔后经咽喉部 、 气管 、 左右支气管进入左右肺叶 ，进入肺叶后支气管进一步细分为更细的支气管 ， 而后进入终末细支气管 。
      \item 从终末细支气管后继续分支 ， 每次分支为两条导管 ， 最后至肺泡管并与肺泡相连 。
      \item 正常人的肺部表面积为 \SI{100}{m^2}
      \item 肺泡是血液与气体进行交换的部位
      \item 巨大的肺泡表面积、丰富的毛细血管和极小的转运距离，决定了肺部给药迅速吸收，且吸收后的药物直接进入血液循环，不受肝首过效应影响。
    \end{itemize}
  \end{minipage}
\end{frame}


\begin{frame}
    \frametitle{在眼部吸收}
    眼部给药主要用于眼局部疾病的治疗 ， 如抗眼部细菌性或病毒性感染 、 降低眼压 、 缩瞳或扩瞳等 。 眼部给药后药物能够到达眼内病灶部位 ， 发挥疾病的治疗作用 。眼睛是人的重要器官 ， 又非常敏感 ， 因此一般不会作为全身治疗作用的给药途径 。
    \hspace{-2em}
    \begin{minipage}{0.38\linewidth}
      \begin{figure}[htpb]
        \centering
        \includegraphics[width=\linewidth]{眼结构}
      \end{figure}
    \end{minipage}\hfill
    \begin{minipage}{0.58\linewidth}
      药物在眼部吸收分为角膜吸收和非角膜吸收两个途径。
      \begin{enumerate}
        \item 角膜吸收一般是眼局部用药的有效吸收途径，角膜外层上皮细胞是药物的主要转运屏障。
        \item 非角膜吸收是药物进入血液循环的主要途径。
      \end{enumerate}
    \end{minipage}
\end{frame}


\begin{frame}
    \frametitle{在眼部吸收}
    \hspace{-2em}
    \begin{minipage}{0.48\linewidth}
      药物在眼部吸收有以下特点：
      \begin{enumerate}
        \item 简单、经济
        \item 可以避免肝首过效应
        \item 对免疫反应不敏感
      \end{enumerate}

      ~~
    \end{minipage}\hfill
    \begin{minipage}{0.48\linewidth}
      眼部给药存在以下缺陷：
      \begin{enumerate}
        \item 眼部刺激性
        \item 剂量损失
        \item 药物在眼部的停留时间问题
        \item 一般病人难以接受眼部给药
      \end{enumerate}
    \end{minipage}
\end{frame}


\begin{frame}
    \frametitle{在鼻腔粘膜中吸收}
    鼻腔给药不仅适用于鼻腔局部疾病的治疗 ， 也是全身疾病治疗的给药途径之一 。 鼻腔给药的药物吸收是药物透过鼻黏膜向循环系统的转运过程。 鼻腔给药的主要优点有 ：

    \begin{enumerate}
      \item 鼻黏膜薄 、 有效表面积大 、 渗透性高 、 血管丰富 ， 有利于吸收 ， 吸收程度和速度有时可与静脉注射相当 ；
      \item 药物直接进入血液循环，可避开药物在胃肠液中的降解和肝首过作用 ；
      \item 能够增加一些药物向脑内传递 ， 有利于脑部疾病的治疗；
      \item 鼻腔内给药方便易行
    \end{enumerate}

    但鼻腔给药也存在不足 ， 如单次用剂量有限 ， 某些药物生物利用度相对较低 ， 吸收剂量不够准确 ， 药物与吸收部位的接触时间相对较短等 。
\end{frame}


\begin{frame}
    \frametitle{在鼻腔粘膜中吸收}
    口服给药个体差异大 、 生物利用度低的药物以及口服易破坏或不吸收 、 只能注射给药的药物 ， 可考虑鼻黏膜给药 。 如在胃肠道中难吸收的磺苄西林 、 头孢唑林可经鼻黏膜吸收 ， 维生素
    \begin{blankblock}{}
      在胃肠道中难吸收的磺苄西林 、 头孢唑林可经鼻黏膜吸收 ， 维生素B12 的鼻用凝胶剂比同剂量口服片剂血药浓度高 $8\sim 10$ 倍 。 药物可制成溶液剂滴入鼻腔 ， 也可以气雾剂给药 。 

      用来治疗头痛和偏头痛的酒石酸布托啡诺鼻喷剂可达到与静脉注射相当的血药浓度水平。
    \end{blankblock}
\end{frame}


\begin{frame}
  \frametitle{在肌肉、皮下吸收}
  \hspace{-2em}
  \begin{minipage}{0.3\linewidth}
    \begin{figure}[htpb]
      \centering
      \includegraphics[width=\linewidth]{注射给药}
      \caption{注射给药示意图}
    \end{figure}
  \end{minipage}\hfill
  \begin{minipage}{0.68\linewidth}
    \begin{enumerate}
      \item 静脉注射药物直接注入血液循环 ， 不存在吸收过程 。静脉注射的生物利用度一般被认为是 100\% 。
      \item 肌内注射存在吸收过程 ， 药物先经注射部位的结缔组织扩散 ， 再经毛细血管吸收进入血液循环 ， 所以肌内注射药物的起效比静脉注射稍慢 ， 但比静脉注射简便安全 ， 比皮下注射刺激性小 ， 应用较广 。肌肉组织内的血管十分丰富 ，一般认为肌内注射给药的药物吸收程度与静脉注射相当 。
    \end{enumerate}
  \end{minipage}

  \begin{enumerate}
    \item[3.] 皮下注射是将药物注射到疏松的皮下组织中 。皮下组织血管较少 ， 血流速度比肌肉组织慢 。 因此 ， 皮下注射后药物吸收较肌内注射慢 ， 有时甚至比口服吸收还慢。皮下植入给药往往可维持数周或数月的疗效。 需延长作用时间的药物可采用皮下注射 ， 如治疗糖尿病的胰岛素等 。
  \end{enumerate}
\end{frame}


\begin{frame}[t]
    \frametitle{}
    \begin{question}
        不存在吸收过程的给药途径是
        \tcblower

        A. 口服给药

        B. 肌内注射

		    C. 肺部给药

        \only<1>{
            D. 静脉注射
        }
        \only<2>{
            \highlight{D. 静脉注射}
        }
    \end{question}
\end{frame}